\section{Task definition and cohort}
\label{sec:task}

\subsection{Outcome: treatment feasibility}
We define \textbf{treatment feasibility} as a binary label $y \in \{0,1\}$ derived from MIMIC-IV v3.1~\citep{johnson2023mimic} discharge and readmission logic (thesis and repository code): $y{=}1$ if the patient is discharged to home, home health care, or rehabilitation within 30 days of ICU admission \emph{and} has no all-cause readmission within 30 days; otherwise $y{=}0$. This operationalization aligns with value-based care incentives but simplifies functional status and social determinants.

\subsection{Inclusion criteria}
We include first ICU stays with: length of stay $\geq 4$\,h, age $\geq 18$, and mappable primary ICD-10 diagnosis for stratification. From eligible stays we draw $n{=}2{,}000$ patients with \textbf{stratified sampling} across five diagnosis groups (sepsis, cardiac, respiratory, renal, diabetes) to preserve rare groups.

\subsection{Features and time window}
For each stay we aggregate \textbf{15 laboratory} and \textbf{7 vital} channels from the first 24\,h of ICU time. Values are clipped to clinician-defined ranges, then min--max normalized to $[0,1]$. We add three engineered features: shock index (HR/SBP), a four-domain SOFA proxy from normalized creatinine, platelets, SpO$_2$, and lactate, and an abnormal-lab fraction. The \textbf{tabular} vector is $\mathbf{x} \in \mathbb{R}^{25}$.

\subsection{Data split and leakage control}
Train/validation/test proportions are \textbf{64\%/16\%/20\%} (1{,}280 / 320 / 400 patients) with stratification by label and node. Random seed is fixed (\texttt{seed=42}). Threshold tuning for classical models uses \emph{only} the validation set; test metrics are reported once. Few-shot models use episodic sampling on the training set and episode-based validation for early stopping.

\subsection{Ethics}
MIMIC-IV is de-identified; use requires PhysioNet credentialing and data use agreement. This secondary analysis does not involve live patient contact.

\begin{figure}[htbp]
  \centering
  \resizebox{\textwidth}{!}{% MIMIC-IV extraction schema — monochrome, publication-style (TikZ)
% Preamble (once per document): \usepackage{tikz}\usetikzlibrary{positioning,arrows.meta,calc}
\begin{tikzpicture}[
  scale=0.88,
  transform shape,
  font=\footnotesize,
  >=Stealth,
  tbl/.style={
    rectangle,
    draw=black!55,
    line width=0.4pt,
    fill=black!3,
    minimum height=8mm,
    inner xsep=4pt,
    inner ysep=3pt,
    align=center,
    rounded corners=1.5pt,
  },
  code/.style={tbl, font=\footnotesize\ttfamily},
  proc/.style={
    rectangle,
    draw=black!55,
    line width=0.45pt,
    fill=black!6,
    inner sep=6pt,
    align=center,
    rounded corners=2pt,
    font=\footnotesize,
  },
  outbox/.style={
    rectangle,
    draw=black!55,
    line width=0.45pt,
    fill=black!10,
    inner sep=6pt,
    align=center,
    rounded corners=2pt,
    font=\footnotesize\bfseries,
  },
  ed/.style={black!50, line width=0.35pt},
  el/.style={font=\footnotesize, text=black!65, inner sep=1pt, fill=white},
]

% --- Core hospital tables ---
\node[code] (pat) {patients};
\node[code, right=9mm of pat] (adm) {admissions};
\node[code, right=9mm of adm] (icu) {icustays};
\node[code, above=9mm of icu, xshift=-6mm] (dx) {diagnoses\_icd};

\draw[ed, ->] (pat) -- node[el, above=0.5pt] {\texttt{subject\_id}} (adm);
\draw[ed, ->] (adm) -- node[el, above=0.5pt] {\texttt{hadm\_id}} (icu);
\draw[ed, ->] (adm.north) |- node[el, pos=0.28, left=1pt] {\texttt{hadm\_id}} (dx.west);

% --- Labs / chart (second row) ---
\node[code, below left=12mm and 2mm of icu] (dlab) {d\_labitems};
\node[code, right=5mm of dlab] (lab) {labevents};
\node[code, right=7mm of lab] (chart) {chartevents};
\node[code, right=5mm of chart] (dit) {d\_items};

\draw[ed, ->] (dlab) -- node[el, above=0.5pt] {\texttt{itemid}} (lab);
\draw[ed, ->] (dit) -- node[el, above=0.5pt] {\texttt{itemid}} (chart);

\coordinate (hub) at ($(icu.south)+(0,-2.5mm)$);
\draw[ed, ->] (icu.south) -- (hub);
\draw[ed, ->] (hub) -| node[el, pos=0.18, below=1pt] {\texttt{stay\_id}} (lab.north);
\draw[ed, ->] (hub) -| (chart.north);

% --- Preprocessing ---
\node[proc, below=11mm of lab, anchor=north, xshift=16mm] (prep) {%
  \textbf{Preprocessing pipeline}\\[0.35em]
  {\footnotesize 15 labs + 7 vitals $\rightarrow$ clip / min--max $[0,1]$ $\rightarrow$ feature engineering}\\[0.2em]
  {\footnotesize $\rightarrow$ \textbf{25-D tabular} $\;|\;$ \textbf{token sequence} $\rightarrow$ ETHOS $\rightarrow$ \textbf{128-D}}%
};

\coordinate (merge) at ($(prep.north)+(0,4.5mm)$);
\draw[ed, ->] (lab.south) |- (merge);
\draw[ed, ->] (chart.south) |- (merge);
\draw[ed, ->] (merge) -- (prep.north);

% --- Cohort output ---
\node[outbox, below=6mm of prep] (cohort) {$2{,}000$ patients $\times$ $25$ features + labels};
\draw[ed, ->] (prep) -- (cohort);

\end{tikzpicture}
}
  \caption{MIMIC-IV extraction and preprocessing schema (TikZ; publication-style monochrome). Join keys and dual outputs match the analysis pipeline.}
  \label{fig:schema}
\end{figure}
